\subsection{Cobertura del corpus, tendencia académica y actividad patentaria}
\label{subsec:results-review}
\zlabel{budget:review-start}

El mapa de cobertura que producen las tres revisiones decide qué comparadores son
legítimos en SP1--SP3, antes de cualquier resultado experimental. La cobertura se mide
por intersecciones, no por capacidad aislada, porque el problema focal es una
composición. El Anexo~\ref{app:review-results-full} discute cada figura panel a panel.

% Figura 3 del consolidado académico: (a) cobertura del núcleo CORE y
% (b) rúbrica de dependencia global. El panel (a) se porta desde
% `academic_literature_review (1).tex` (líneas 228-238); el panel (b) es nuevo
% en el PDF vigente y se reconstruye aquí con sus tres categorías literales.
\begin{figure}[!htbp]
  \centering
  \begin{minipage}[t]{0.46\linewidth}
    \centering
    {\sffamily\bfseries\fontsize{8.4}{9}\selectfont (a) Cobertura del corpus CORE\par}
    \vspace{2mm}
    \resizebox{\linewidth}{!}{%
\begin{tikzpicture}[x=1.72mm,y=6.45mm,font=\sffamily\fontsize{5.95}{6.35}\selectfont]
\foreach \lab/\v/\y in {Transporte compartido/23/7,Reclutamiento/18/6,Certificación física/13/5,Compartido + distribuido/10/4,Heterog. + coalición variable/9/3,Compartido + reemplazo/2/2,Heterog. + compartido/0/1,Física + distrib. + reemplazo/0/0}{
  \fill[onyx!6] (0,\y-.30) rectangle (24.8,\y+.30);
  \fill[orange] (0,\y-.30) rectangle (\v,\y+.30);
  \node[anchor=east,align=right,text width=31mm] at (-1.0,\y) {\lab};
  \node[anchor=west,font=\bfseries] at (\v+.55,\y) {\v};}
\draw[-{Stealth[length=1.45mm]},ink,line width=.4pt](0,-.65)--(25.3,-.65);
\node[anchor=north]at(12.2,-.82){trabajos CORE};
\end{tikzpicture}}
  \end{minipage}\hfill
  \begin{minipage}[t]{0.50\linewidth}
    \centering
    {\sffamily\bfseries\fontsize{8.4}{9}\selectfont (b) Rúbrica de dependencia global\par}
    \vspace{2mm}
    \resizebox{\linewidth}{!}{%
\begin{tikzpicture}[x=1mm,y=1mm,font=\sffamily\fontsize{6.0}{6.6}\selectfont]
\path[use as bounding box] (0,-8) rectangle (100,46);
\foreach \y/\bg/\fg/\name/\txt in {%
  34/{softblue}/{gamecol}/{LOCAL}/{decisión y cierre con estado propio/vecinal; sin servicio global permanente},
  19/{soft}/{muted}/{MIXTA}/{acción local, pero con líder, precondición, mapa o cierre global material},
  4/{softwarm}/{orange}/{GLOBAL}/{asignación, replanificación o selección decididos por una entidad central en tiempo de ejecución}}{%
  \fill[\bg,rounded corners=1pt] (0,\y) rectangle (100,\y+11);
  \node[anchor=west,font=\bfseries\fontsize{9}{9}\selectfont,text=\fg] at (3,\y+5.5) {\name};
  \node[anchor=west,align=left,text width=62mm,text=ink] at (26,\y+5.5) {\txt};}
\node[anchor=north,align=center,text width=86mm,text=muted,font=\fontsize{5.4}{6.0}\selectfont] at (50,1)
  {La categoría se asigna independientemente a Recruit, Certify, Dock, Transport y Recover; no se promedia entre etapas.};
\end{tikzpicture}}
  \end{minipage}
  \caption[Síntesis estructural del núcleo curado.]{Síntesis estructural del
  núcleo curado. (a) Cobertura de capacidades e intersecciones: el corpus cubre
  con solvencia cada capacidad por separado y la cobertura se desploma al
  exigirlas simultáneamente. (b) Rúbrica usada para codificar la dependencia
  global por etapa. La categoría se asigna por trabajo y por etapa; no se
  promedian categorías ordinales ni se interpreta su separación como distancia
  métrica.}
  \label{fig:lit-coverage-rubric}
  \viusource{elaboración propia a partir del dataset académico y de la
  codificación del núcleo canónico}
\end{figure}


La cobertura por capacidad aislada es solvente (23 trabajos en transporte compartido, 18
en reclutamiento) y colapsa a dos al exigir carga compartida y reemplazo, y a cero al
sumar certificación física (Figura~\ref{fig:lit-coverage-rubric}).

% Figura 6 del consolidado académico (PDF vigente, página 6): cobertura de
% etapas, intersecciones funcionales y señales positivas de evidencia.
% Reconstruida en LaTeX a partir de los valores impresos en el original.
\begin{figure}[!htbp]
  \centering
  \begin{minipage}[t]{0.47\linewidth}
    \centering
    {\sffamily\bfseries\fontsize{7.6}{8.2}\selectfont (a) SP1 concentra la cobertura\par}
    \vspace{1.5mm}
    \resizebox{\linewidth}{!}{%
\begin{tikzpicture}[x=1mm,y=1mm,font=\sffamily\fontsize{6.0}{6.5}\selectfont]
% La escala deja sitio al rótulo del valor más alto dentro del cuadro:
% 115 documentos ocupan 55 mm y la etiqueta cabe antes del borde en 100 mm.
\path[use as bounding box] (0,-9) rectangle (100,34);
\def\sc{0.478}
\foreach \y/\lab/\v/\col in {26/{SP1 · coalición / asignación}/115/{gamecol},14/{SP3 · planificación / tráfico}/20/{canary!80!onyx},2/{SP2 · transporte físico}/15/{autumn}}{%
  \fill[onyx!5] (36,\y-3.4) rectangle (36+\sc*122,\y+3.4);
  \fill[\col,rounded corners=.5pt] (36,\y-3.4) rectangle (36+\sc*\v,\y+3.4);
  \node[anchor=east,text=ink] at (34.5,\y) {\lab};
  \node[anchor=west,font=\bfseries,text=ink] at (36+\sc*\v+1.4,\y) {\v};}
\draw[-{Stealth[length=1.4mm]},ink,line width=.35pt] (36,-3)--(36+\sc*124,-3);
\node[anchor=north,text=muted,font=\fontsize{5.2}{5.6}\selectfont] at (36+\sc*62,-4){documentos del corpus};
\end{tikzpicture}}
  \end{minipage}\hfill
  \begin{minipage}[t]{0.50\linewidth}
    \centering
    {\sffamily\bfseries\fontsize{7.6}{8.2}\selectfont (b) La integración entre etapas es excepcional\par}
    \vspace{1.5mm}
    \resizebox{\linewidth}{!}{%
\begin{tikzpicture}[x=1mm,y=1mm,font=\sffamily\fontsize{5.8}{6.2}\selectfont]
% La escala vertical deja 6 mm libres sobre la barra de 103 para su rótulo y
% para la nota, de modo que el panel no invade al de la izquierda.
\path[use as bounding box] (-6,-14) rectangle (104,40);
\def\sy{0.245}
\foreach \i/\lab/\v in {0/{solo\\SP1}/103,1/{solo\\SP2}/13,2/{solo\\SP3}/10,3/{SP1 +\\SP3}/10,4/{SP1 +\\SP2}/2,5/{SP2 +\\SP3}/0,6/{las\\3}/0}{%
  \pgfmathsetmacro{\x}{10+\i*13}
  \pgfmathsetmacro{\h}{\sy*\v}
  \ifdim \v pt>0pt
    \fill[onyx!42,rounded corners=.5pt] (\x,0) rectangle (\x+9,\h);
  \fi
  \node[anchor=south,font=\bfseries\fontsize{5.8}{6.2}\selectfont,text=ink] at (\x+4.5,\h+0.8) {\v};
  \node[anchor=north,align=center,text=muted,font=\fontsize{5.0}{5.4}\selectfont] at (\x+4.5,-6.2) {\lab};}
% marcadores de etapa bajo cada columna
\foreach \i/\a/\b/\c in {0/1/0/0,1/0/1/0,2/0/0/1,3/1/0/1,4/1/1/0,5/0/1/1,6/1/1/1}{%
  \pgfmathsetmacro{\x}{10+\i*13}
  \foreach \k/\col/\on in {0/{gamecol}/\a,1/{autumn}/\b,2/{canary!80!onyx}/\c}{%
    \pgfmathsetmacro{\yy}{-1.6-\k*1.9}
    \ifnum\on=1 \fill[\col] (\x+4.5,\yy) circle (0.75);
    \else \draw[onyx!35,line width=.22pt] (\x+4.5,\yy) circle (0.75);\fi}}
\node[anchor=east,text=muted,font=\fontsize{4.6}{5.0}\selectfont] at (8,-1.6){SP1};
\node[anchor=east,text=muted,font=\fontsize{4.6}{5.0}\selectfont] at (8,-3.5){SP2};
\node[anchor=east,text=muted,font=\fontsize{4.6}{5.0}\selectfont] at (8,-5.4){SP3};
\node[anchor=north west,align=left,text width=62mm,text=muted,font=\fontsize{5.0}{5.6}\selectfont] at (30,38)
  {103 trabajos cubren solo SP1; solo 2 enlazan SP1--SP2 y ninguno cubre las tres etapas.};
\node[rotate=90,anchor=south,text=muted,font=\fontsize{5.2}{5.6}\selectfont] at (-4,14){documentos};
\end{tikzpicture}}
  \end{minipage}

  \vspace{2.5mm}
  {\sffamily\bfseries\fontsize{7.6}{8.2}\selectfont (c) La evidencia física y formal aparece con mucha menor frecuencia\par}
  \vspace{1.5mm}
  \makebox[\textwidth][c]{\resizebox{0.86\textwidth}{!}{%
\begin{tikzpicture}[x=1mm,y=1mm,font=\sffamily\fontsize{6.0}{6.5}\selectfont]
\def\sc{1.30}
\foreach \y/\lab/\v/\col in {24/{Simulación}/99/{learncol},16/{Contexto / caso industrial}/25/{mustard!85!onyx},8/{Experimento físico}/12/{autumn},0/{Análisis formal}/7/{onyx!45}}{%
  \fill[onyx!5] (40,\y-2.6) rectangle (40+\sc*105,\y+2.6);
  \fill[\col,rounded corners=.5pt] (40,\y-2.6) rectangle (40+\sc*\v,\y+2.6);
  \node[anchor=east,text=ink] at (38.5,\y) {\lab};
  \node[anchor=west,font=\bfseries,text=ink] at (40+\sc*\v+1.4,\y) {\v};}
\foreach \v in {0,20,40,60,80,100}{%
  \node[anchor=north,text=muted,font=\fontsize{5.2}{5.6}\selectfont] at (40+\sc*\v,-3.4){\v};}
\node[anchor=north west,text=muted,font=\fontsize{5.0}{5.4}\selectfont] at (40,-7)
  {«Contexto / caso industrial» puede acompañar a simulación; son conteos positivos, no prevalencias.};
\end{tikzpicture}%
  }}
  \caption[Cobertura de etapas, intersecciones y señales de evidencia.]{Cobertura
  de etapas, intersecciones funcionales y señales positivas de evidencia en el
  corpus analítico. (a) Cobertura codificada por etapa. (b) Intersecciones
  observadas bajo la taxonomía SP1--SP3; los ceros significan «no identificado
  bajo esta codificación», no inexistencia bibliográfica. (c) Conteos positivos
  por tipo de evidencia; no se presentan porcentajes de prevalencia porque la
  disponibilidad de extracción estructural varía por campo.}
  \label{fig:lit-stage-coverage}
  \viusource{tablas derivadas del mapeo analítico, $n=244$}
\end{figure}


La cobertura se concentra en SP1 (115 señales frente a 20 de SP3 y 15 de SP2); ningún
documento cubre las tres etapas y la validación es sobre todo por simulación, 99 señales
frente a 12 de experimento físico (Figura~\ref{fig:lit-stage-coverage}).

% Figura 5 del consolidado académico (PDF vigente, página 6): cambio
% metodológico por periodos. Reconstruida en LaTeX a partir de los valores
% impresos en el original. Los cuatro paneles comparten escala 0--35 % para que
% la comparación entre periodos sea directa.
%
% Cada fila se dibuja con \mcrow, que recibe el índice de fila, la etiqueta, los
% cuatro valores y el cambio. No se anida \foreach dentro de \foreach: PGF no
% expande de forma fiable los valores del bucle exterior dentro de la lista del
% interior, y el resultado era un «You can't use the character» por fila.
\begin{figure}[!htbp]
  \centering
  {\sffamily\bfseries\fontsize{8.6}{9.2}\selectfont Cambio metodológico por periodos\par}
  {\sffamily\fontsize{6.2}{6.8}\selectfont\color{muted}Misma escala 0--35~\% en los cuatro periodos; la columna final muestra el cambio entre el primero y el último.\par}
  \vspace{1.5mm}
  \makebox[\textwidth][c]{\resizebox{\textwidth}{!}{%
\begin{tikzpicture}[x=1mm,y=1mm,font=\sffamily\fontsize{5.6}{6.0}\selectfont]
  \def\panw{33}\def\gap{5}\def\rowh{6.0}\def\sc{0.943}
  \def\xzero{169}% origen del panel de cambio: 4*(33+5)+17

  % Una barra: #1 origen x del panel, #2 valor, #3 color.
  \def\mcbar#1#2#3#4{%
    \pgfmathsetmacro{\bw}{\sc*#2}%
    \ifdim #2pt>0pt
      \fill[#3,rounded corners=.4pt] (#1,\yy-2.0) rectangle (#1+\bw,\yy+2.0);
      \ifdim \bw pt>13pt
        \node[anchor=east,font=\bfseries\fontsize{5.6}{6.0}\selectfont,text=#4] at (#1+\bw-1.2,\yy){#2};
      \else
        \node[anchor=west,font=\fontsize{5.6}{6.0}\selectfont,text=ink] at (#1+\bw+1.2,\yy){#2};
      \fi
    \fi}

  % Una fila completa: #1 índice, #2 etiqueta, #3-#6 valores, #7 cambio.
  \def\mcrow#1#2#3#4#5#6#7{%
    \pgfmathsetmacro{\yy}{54-#1*\rowh}%
    \node[anchor=east,text=ink,font=\fontsize{6.0}{6.4}\selectfont] at (-1.5,\yy){#2};
    \mcbar{0}{#3}{gamecol!22}{ink}%
    \mcbar{38}{#4}{gamecol!45}{ink}%
    \mcbar{76}{#5}{gamecol!85}{white}%
    \mcbar{114}{#6}{autumn}{white}%
    \pgfmathsetmacro{\dw}{1.1*#7}%
    \ifdim #7pt>0pt
      \fill[canary!80!onyx,rounded corners=.4pt] (\xzero,\yy-2.0) rectangle (\xzero+\dw,\yy+2.0);
      \node[anchor=west,font=\bfseries\fontsize{5.8}{6.2}\selectfont,text=ink] at (\xzero+\dw+1.2,\yy){+#7};
    \else
      \fill[autumn,rounded corners=.4pt] (\xzero+\dw,\yy-2.0) rectangle (\xzero,\yy+2.0);
      \node[anchor=east,font=\bfseries\fontsize{5.8}{6.2}\selectfont,text=ink] at (\xzero+\dw-1.2,\yy){#7};
    \fi}

  % Fondos, rejilla y cabeceras de los cuatro periodos.
  \foreach \p/\lab/\n in {0/{Hasta 2010}/{$n=35$},1/{2011--2016}/{$n=40$},2/{2017--2021}/{$n=89$},3/{2022--2026}/{$n=81$}}{%
    \pgfmathsetmacro{\px}{\p*(\panw+\gap)}
    \node[anchor=south,font=\bfseries\fontsize{6.4}{6.8}\selectfont,text=ink] at (\px+\panw/2,59.4){\lab};
    \node[anchor=south,font=\fontsize{5.0}{5.4}\selectfont,text=muted] at (\px+\panw/2,56.6){\n};
    \fill[onyx!3] (\px,-1) rectangle (\px+\panw,56);
    \foreach \v in {0,10,20,30}{%
      \draw[rule!60,line width=.2pt] (\px+\sc*\v,-1)--(\px+\sc*\v,56);
      \node[anchor=north,text=muted,font=\fontsize{5.0}{5.4}\selectfont] at (\px+\sc*\v,-1.6){\v};}}
  \node[anchor=south,font=\bfseries\fontsize{6.4}{6.8}\selectfont,text=ink] at (\xzero,59.4){Cambio temprano--reciente};
  \node[anchor=south,font=\fontsize{5.0}{5.4}\selectfont,text=muted] at (\xzero,56.6){puntos porcentuales};

  \mcrow{0}{Consenso / distribución}{20.0}{20.0}{31.6}{24.7}{4.7}
  \mcrow{1}{Subasta / mercado}{22.9}{17.5}{16.5}{13.5}{-9.4}
  \mcrow{2}{Aprendizaje}{2.9}{5.0}{17.7}{15.7}{12.8}
  \mcrow{3}{Metaheurística evolutiva}{0.0}{5.0}{16.5}{11.2}{11.2}
  \mcrow{4}{Control de formación}{14.3}{5.0}{5.1}{7.9}{-6.4}
  \mcrow{5}{Comportamiento / enjambre}{8.6}{5.0}{0.0}{5.6}{-3.0}
  \mcrow{6}{Teoría de juegos}{2.9}{5.0}{3.8}{3.4}{0.5}
  \mcrow{7}{Seguridad reactiva / CBF}{0.0}{2.5}{2.5}{6.7}{6.7}
  \mcrow{8}{Optimización exacta}{5.7}{0.0}{5.1}{2.2}{-3.5}

  \draw[onyx,line width=.45pt] (\xzero,-1)--(\xzero,56);
  \foreach \v in {-10,0,10}{%
    \node[anchor=north,text=muted,font=\fontsize{5.0}{5.4}\selectfont] at (\xzero+1.1*\v,-1.6){\v};}
\end{tikzpicture}%
  }}
  \caption[Cambio metodológico por periodos.]{Cambio metodológico por periodos.
  Los cuatro paneles temporales comparten la misma escala horizontal, expresada
  como porcentaje interno de cada periodo, para comparar directamente la
  persistencia y la recombinación de familias; el panel final muestra el cambio
  en puntos porcentuales entre el periodo temprano y 2022--2026. Las cuotas son
  proporciones dentro de cada periodo, no tasas de publicación, y la
  clasificación es multietiqueta.}
  \label{fig:lit-method-change}
  \viusource{elaboración propia a partir del corpus analítico; 2026 es un año incompleto}
\end{figure}


% Sustituto real de los paneles (b)/(c) de la Figura 4 original.
%
% El PDF `MROB_literature_review_reworked (2).pdf` usa una taxonomía de siete
% familias (Coalition, Transport, Planning, Feasibility, Recovery, Learning/RL,
% Foundational) que no tiene dataset reproducible en el repositorio: la propia
% auditoría del proyecto (academic-review/literature-review-v4/README.md)
% confirma que esas cifras no se derivaron de ningún CSV conservado. Esta
% figura usa en su lugar las cinco categorías reales del corpus V3
% (`analytic_corpus_v3.csv`, campo `theme_signals_title_abstract`), generadas
% por `pre-thesis/scripts/build_review_figure_data.py` y volcadas en
% `shared/generated-macros/lit-corpus-activity.tex`. La clasificación es
% multietiqueta, así que las cuotas de un mismo periodo pueden sumar más de
% 100~%.
\begin{figure}[!htbp]
  \centering
  \begin{minipage}[t]{0.56\linewidth}
    \centering
    {\sffamily\bfseries\fontsize{7.8}{8.4}\selectfont (a) Cuota temática por periodo\par}
    \vspace{1.5mm}
    \resizebox{\linewidth}{!}{%
\begin{tikzpicture}[x=1mm,y=1mm,font=\sffamily\fontsize{5.8}{6.2}\selectfont]
  \def\sc{0.62}
  \foreach \p/\lab in {0/{$\le$2010}, 1/{2011--16}, 2/{2017--21}, 3/{2022--26}}{%
    \pgfmathsetmacro{\y}{34-\p*9}
    \node[anchor=east,text=ink,font=\bfseries] at (-2,\y+3.4){\lab};}
  \foreach \y/\lab/\col/\vals in {%
    34/{Coalición / MRTA}/{gamecol}/{45.7,50.0,48.1,44.9},
    24.4/{Recuperación / resil.}/{canary!75!onyx}/{40.0,25.0,35.4,30.3},
    14.8/{Contexto / otros}/{VIUGray!70}/{17.1,40.0,24.1,34.8},
    5.2/{Planif. / rutas}/{VIUOrange}/{2.9,7.5,11.4,7.9},
    -4.4/{Transporte / control}/{SPPurple}/{8.6,2.5,5.1,7.9}}{%
    \node[anchor=east,text=ink] at (-2,\y){\lab};
    \foreach \v [count=\i from 0] in \vals{%
      \pgfmathsetmacro{\x}{\i*15}
      \fill[onyx!5] (\x,\y-1.7) rectangle (\x+13,\y+1.7);
      \pgfmathsetmacro{\w}{\sc*\v}
      \fill[\col,rounded corners=.4pt] (\x,\y-1.7) rectangle (\x+\w,\y+1.7);
      \node[anchor=west,font=\fontsize{5.0}{5.4}\selectfont,text=ink] at (\x+\w+.6,\y){\v};}}
  \foreach \p/\lab in {0/{$\le$2010},1/{2011-16},2/{2017-21},3/{2022-26}}{%
    \pgfmathsetmacro{\x}{\p*15}
    \node[anchor=south,font=\fontsize{5.2}{5.6}\selectfont,text=muted] at (\x+6.5,37.4){\lab};}
\end{tikzpicture}}
  \end{minipage}\hfill
  \begin{minipage}[t]{0.40\linewidth}
    \centering
    {\sffamily\bfseries\fontsize{7.8}{8.4}\selectfont (b) Cambio temprano--reciente\par}
    {\sffamily\fontsize{5.6}{6.0}\selectfont\color{muted}$\le$2010 $\rightarrow$ 2022--2026, puntos porcentuales\par}
    \vspace{1.5mm}
    \resizebox{\linewidth}{!}{%
\begin{tikzpicture}[x=1mm,y=1mm,font=\sffamily\fontsize{5.8}{6.2}\selectfont]
% La caja terminaba en x=68 y la etiqueta «+17.7», anclada al oeste en x=63.4,
% quedaba fuera: al reescalar a \linewidth invadia el margen derecho.
\path[use as bounding box] (0,-6) rectangle (74,36);
  \draw[VIUGray!55,line width=.35pt] (40,-4)--(40,34);
  % Cinco filas fijas: la escala y la posición de cada etiqueta se calculan a
  % mano, no con una fórmula genérica; con solo cinco valores conocidos y dos
  % de ellos casi nulos (-0,7 y -0,8) una regla \ifdim/min/max encadenada con
  % signos negativos es más frágil que útil. La etiqueta de categoría se ancla
  % en x=20, lejos de las etiquetas de valor (x=26,5-37,5): a x=36, como en el
  % primer intento, ambas colisionaban y el solape de trazos parecía una
  % fuente gigante, aunque el PDF confirma que las cinco miden 10,55 pt.
  \node[anchor=east,text=ink] at (14,30){Contexto / otros};
  \fill[canary!80!onyx,rounded corners=.4pt] (40,28) rectangle (61.9,32);
  \node[anchor=west,font=\bfseries,text=ink] at (63.4,30){+17.7};

  \node[anchor=east,text=ink] at (14,22){Planif. / rutas};
  \fill[VIUOrange,rounded corners=.4pt] (40,20) rectangle (46.2,24);
  \node[anchor=west,font=\bfseries,text=ink] at (47.7,22){+5.0};

  \node[anchor=east,text=ink] at (14,14){Transporte / control};
  \fill[SPPurple,rounded corners=.4pt] (39.1,12) rectangle (40,16);
  \node[anchor=east,font=\bfseries,text=ink] at (37.5,14){-0.7};

  \node[anchor=east,text=ink] at (14,6){Coalición / MRTA};
  \fill[gamecol,rounded corners=.4pt] (39.0,4) rectangle (40,8);
  \node[anchor=east,font=\bfseries,text=ink] at (37.5,6){-0.8};

  \node[anchor=east,text=ink] at (14,-2){Recuperación / resil.};
  \fill[SPRed,rounded corners=.4pt] (28.0,-4) rectangle (40,0);
  \node[anchor=east,font=\bfseries,text=ink] at (26.5,-2){-9.7};
\end{tikzpicture}}
  \end{minipage}
  \caption[Cuota temática por periodo y cambio temprano-reciente.]{Cuota temática
  por periodo (a) y cambio entre el periodo temprano y 2022--2026 (b), calculados
  sobre las cinco categorías reales del corpus analítico V3. La clasificación es
  multietiqueta: las cuotas de un mismo periodo pueden sumar más de 100~\%. Esta
  figura no reproduce las siete familias del panel visual original, que carecen
  de dataset conservado; usa en su lugar la taxonomía verificable del corte V3.}
  \label{fig:lit-period-shift}
  \viusource{elaboración propia a partir de \texttt{analytic\_corpus\_v3.csv} y
  \texttt{theme\_period\_v3.csv}, generado por
  \texttt{pre-thesis/scripts/build\_review\_figure\_data.py}}
\end{figure}


El cambio metodológico por periodo es de recombinación, no de sustitución: consenso y
distribución, subasta y aprendizaje coexisten con pesos que se redistribuyen, sin que
ninguna familia converja (Figuras~\ref{fig:lit-method-change}
y~\ref{fig:lit-period-shift}).

% Sustituto real de la Figura 8 original (estructura bibliométrica del corpus).
%
% El panel (a) original identifica núcleos A1-A11 mediante una ejecución de
% Louvain con semilla propia que no está guardada en el repositorio: no hay
% forma de reproducirla. En su lugar, el panel (a) de aquí calcula las
% componentes conexas del grafo de coautoría fuerte (par de autores con 3 o
% más documentos compartidos en el corpus V3), que es un criterio explícito y
% reproducible. El panel (b) sí usa el mismo tipo de insumo que el original: la
% coocurrencia de familias de método, con comunidades obtenidas por modularidad
% voraz (Clauset-Newman-Moore) sobre `method_cooccurrence_v3.csv`.
\begin{figure}[!htbp]
  \centering
  \begin{minipage}[t]{0.46\linewidth}
    \centering
    {\sffamily\bfseries\fontsize{7.8}{8.4}\selectfont (a) Núcleos de coautoría fuerte\par}
    {\sffamily\fontsize{5.6}{6.0}\selectfont\color{muted}Pares con \LitCoauthorThreshold{} o más documentos compartidos ($n=$\LitCoauthorPairsTotal{} pares totales)\par}
    \vspace{2mm}
    \resizebox{\linewidth}{!}{%
\begin{tikzpicture}[x=1mm,y=1mm,font=\sffamily\fontsize{5.4}{5.8}\selectfont,
  auth/.style={circle,draw=gamecol,fill=gamecol!12,line width=.5pt,minimum size=5mm,inner sep=0pt,font=\fontsize{4.6}{4.9}\selectfont},
  wedge/.style={draw=gamecol!70,line width=.4pt}]
\path[use as bounding box] (0,-4) rectangle (100,42);
% Núcleo 1: K4 (Lefebvre, Leclercq, Guérin, Chakraa), todas las aristas peso 3.
\node[auth] (l1) at (14,32) {DL};
\node[auth] (l2) at (30,32) {EL};
\node[auth] (l3) at (14,18) {FG};
\node[auth] (l4) at (30,18) {HC};
\foreach \a/\b in {l1/l2,l1/l3,l1/l4,l2/l3,l2/l4,l3/l4}{\draw[wedge] (\a)--(\b);}
\node[anchor=north,align=center,text width=32mm,text=muted] at (22,10)
  {\textbf{4 autores, $K_4$}\\Lefebvre, Leclercq, Guérin, Chakraa\\todas las aristas: 3 documentos};
% Núcleo 2: triángulo (Zitouni, Maamri, Harous), pesos 6/4/4.
\node[auth] (z1) at (58,34) {FZ};
\node[auth] (z2) at (48,18) {RM};
\node[auth] (z3) at (68,18) {SH};
\draw[wedge,line width=.75pt] (z1)--(z2) node[midway,anchor=south east,font=\fontsize{4.2}{4.5}\selectfont,text=gamecol]{6};
\draw[wedge] (z1)--(z3) node[midway,anchor=south west,font=\fontsize{4.2}{4.5}\selectfont,text=gamecol]{4};
\draw[wedge] (z2)--(z3) node[midway,anchor=north,font=\fontsize{4.2}{4.5}\selectfont,text=gamecol]{4};
\node[anchor=north,align=center,text width=30mm,text=muted] at (58,10)
  {\textbf{3 autores, triángulo}\\Zitouni, Maamri, Harous\\6, 4 y 4 documentos};
% Núcleo 3: un par aislado (Shell, Liu).
\node[auth] (s1) at (88,28) {DS};
\node[auth] (s2) at (88,18) {LL};
\draw[wedge] (s1)--(s2) node[midway,anchor=west,font=\fontsize{4.2}{4.5}\selectfont,text=gamecol]{3};
\node[anchor=north,align=center,text width=22mm,text=muted] at (88,10)
  {\textbf{par aislado}\\Shell, Liu\\3 documentos};
\end{tikzpicture}}
  \end{minipage}\hfill
  \begin{minipage}[t]{0.50\linewidth}
    \centering
    {\sffamily\bfseries\fontsize{7.8}{8.4}\selectfont (b) Comunidades método--coocurrencia\par}
    {\sffamily\fontsize{5.6}{6.0}\selectfont\color{muted}Modularidad voraz, $Q=$\LitMethodModularity{} sobre \LitMethodCommunities{} comunidades\par}
    \vspace{2mm}
    \resizebox{\linewidth}{!}{%
\begin{tikzpicture}[x=1mm,y=1mm,font=\sffamily\fontsize{6.0}{6.5}\selectfont]
\node[rounded corners=2pt,draw=gamecol,fill=gamecol!8,line width=.5pt,
  minimum width=52mm,minimum height=26mm,align=left,text width=48mm,inner sep=2mm]
  (ca) at (28,20) {\textbf{\textcolor{gamecol}{Comunidad A}}\\[1mm]
  Consenso · Subasta / mercado\\Metaheurística · Optimización exacta\\Control predictivo};
\node[rounded corners=2pt,draw=controlcol,fill=controlcol!8,line width=.5pt,
  minimum width=52mm,minimum height=26mm,align=left,text width=48mm,inner sep=2mm]
  (cb) at (28,-10) {\textbf{\textcolor{controlcol}{Comunidad B}}\\[1mm]
  Aprendizaje · Enjambre · Formación\\Teoría de juegos\\\emph{Graph search} / MAPF · CBF};
\draw[-{Stealth[length=1.6mm]},VIUGray,line width=.4pt] (ca.south)--(cb.north)
  node[midway,anchor=west,align=left,font=\fontsize{5.0}{5.4}\selectfont,text=muted]{18 documentos comparten\\subasta + consenso};
\end{tikzpicture}}
  \end{minipage}
  \caption[Estructura bibliométrica del corpus: coautoría y método.]{Estructura
  bibliométrica del corpus analítico V3. (a) Núcleos de coautoría con \LitCoauthorThreshold{}
  o más documentos compartidos, calculados como componentes conexas del grafo
  de coautoría fuerte. (b) Comunidades de familias de método obtenidas por
  modularidad voraz sobre su coocurrencia documental; $Q=$\LitMethodModularity{}
  es baja, señal de que las familias metodológicas se combinan con frecuencia y
  no forman bloques aislados. Ninguno de los dos paneles reproduce los
  identificadores A1--A11 del panel visual original: esa partición usó una
  ejecución de Louvain sin semilla ni dataset conservados en el repositorio.}
  \label{fig:lit-bibliometric}
  \viusource{elaboración propia a partir de \texttt{top\_coauthor\_pairs\_v3.csv}
  y \texttt{method\_cooccurrence\_v3.csv}, generado por
  \texttt{pre-thesis/scripts/build\_review\_figure\_data.py}}
\end{figure}


La estructura bibliométrica confirma esa lectura: núcleos de coautoría pequeños y
aislados, y comunidades de método de modularidad baja ($Q=$\LitMethodModularity{}), sin
una agenda ni una escuela metodológica dominante (Figura~\ref{fig:lit-bibliometric}).

% Figura 2 del consolidado industrial, portada desde
% `thesis-reviews/TFM_estado_industrial_patentes_EDITABLE_v47.tex` (líneas
% 351-452): (a) actividad anual con MM3 e inset de cambio temático,
% (b) cobertura geográfica y (c) orientación funcional por dominio.
% Requiere pgfplots y la librería patterns.meta, ambas cargadas en main-v2.tex.
\begin{figure}[!htbp]
  \centering
{\sffamily\fontsize{8.45}{9.0}\selectfont\bfseries (a) Publicaciones anuales y cambio del foco temático del corpus}\par\vspace{1pt}
\makebox[\textwidth][c]{\resizebox{\textwidth}{!}{%
\begin{tikzpicture}
\begin{axis}[
 name=patents,width=0.95\linewidth,height=6.12cm,
 ymin=0,ymax=31,ylabel={\sffamily\fontsize{6.8}{7.2}\selectfont Publicaciones del corpus/año},
 symbolic x coords={2012,2013,2014,2015,2016,2017,2018,2019,2020,2021,2022,2023,2024,2025},xtick=data,
 axis x line*=bottom,axis y line*=left,
 x tick label style={font=\sffamily\fontsize{5.6}{5.95}\selectfont,rotate=45,anchor=east,text=muted},
 y tick label style={font=\sffamily\fontsize{5.8}{6.1}\selectfont,text=muted},ytick={0,5,10,15,20,25,30},
 axis line style={draw=rule!95,line width=.35pt},tick style={draw=rule!95},
 ymajorgrids=true,major grid style={draw=rule!50,line width=.18pt},enlarge x limits=0.025,clip=false,
 legend style={draw=none,fill=none,font=\sffamily\fontsize{5.25}{5.55}\selectfont,at={(0.52,-0.16)},anchor=north,legend columns=3,/tikz/every even column/.append style={column sep=1.2mm}}]
% ---- Annual macro-layer counts; edit coordinates here ----
\addplot+[ybar stacked,bar width=10.0pt,fill=autumn,draw=white,line width=.22pt] coordinates {(2012,0)(2013,0)(2014,3)(2015,0)(2016,0)(2017,5)(2018,3)(2019,4)(2020,8)(2021,6)(2022,11)(2023,11)(2024,12)(2025,17)};
\addlegendentry{Decisión de misión}
\addplot+[ybar stacked,bar width=10.0pt,fill=mustard,draw=white,line width=.22pt] coordinates {(2012,1)(2013,0)(2014,0)(2015,0)(2016,1)(2017,1)(2018,0)(2019,3)(2020,1)(2021,8)(2022,11)(2023,5)(2024,4)(2025,3)};
\addlegendentry{Coordinación cinemática}
\addplot+[ybar stacked,bar width=10.0pt,fill=onyx!42!white,draw=white,line width=.22pt] coordinates {(2012,1)(2013,0)(2014,0)(2015,2)(2016,1)(2017,1)(2018,3)(2019,4)(2020,10)(2021,7)(2022,4)(2023,7)(2024,7)(2025,2)};
\addlegendentry{Interacción física}
% MM3 volume, left axis
\addplot+[no markers,smooth,onyx,line width=.9pt,dash pattern=on 3pt off 2pt] coordinates {(2014,1.67)(2015,1.67)(2016,2.33)(2017,3.67)(2018,5.00)(2019,8.00)(2020,12.00)(2021,17.00)(2022,22.00)(2023,23.33)(2024,24.00)(2025,22.67)};
\node[anchor=west,font=\sffamily\fontsize{5.0}{5.3}\selectfont,text=onyx] at (axis description cs:0.70,0.94){MM3 volumen};
% Total labels
\addplot+[draw=none,forget plot,nodes near coords,point meta=explicit symbolic,every node near coord/.append style={font=\sffamily\fontsize{5.0}{5.3}\selectfont\bfseries,text=charcoal,yshift=1.2pt}] coordinates {(2012,2)[2](2013,0)[ ](2014,3)[3](2015,2)[2](2016,2)[2](2017,7)[7](2018,6)[6](2019,11)[11](2020,19)[19](2021,21)[21](2022,26)[26](2023,23)[23](2024,23)[23](2025,22)[22]};
% Inset: six thematic shifts, p.p. (fixed columns: topic | bar | delta)
\node[anchor=north west,inner sep=0pt] at (axis description cs:-0.085,0.97){%
\begin{tikzpicture}[x=1mm,y=1mm,font=\sffamily]
 \path[use as bounding box] (0,1.2) rectangle (61,31);
 \node[anchor=west,font=\fontsize{5.65}{6}\selectfont\bfseries,text=spicy] at (15.0,29.1){Cambio de cuota temática [p.p.]};
 \node[anchor=west,font=\fontsize{4.45}{4.75}\selectfont,text=muted] at (15.0,26.2){2020--2022 $\rightarrow$ 2023--2025};
 \node[anchor=east,font=\fontsize{4.0}{4.3}\selectfont\bfseries,text=muted] at (59.0,24.6){$\Delta$ p.p.};
 \draw[onyx!32,line width=.28pt] (36,2.0)--(36,24.0);
 \foreach \y/\lab in {22.4/{Planificación},18.7/{Factibilidad},15.0/{Reclutamiento},11.3/{Acoplamiento},7.6/{Formación},3.9/{Shared-load}}{\node[anchor=east,font=\fontsize{4.4}{4.7}\selectfont,text=ink] at (27.8,\y){\lab};}
 \fill[autumn] (36,21.65) rectangle ({36+0.55*16.67},23.15); \node[anchor=east,font=\fontsize{4.3}{4.6}\selectfont,text=autumn] at (59.0,22.4){+16,7};
 \fill[canary] (36,17.95) rectangle ({36+0.55*5.84},19.45); \node[anchor=east,font=\fontsize{4.3}{4.6}\selectfont,text=green!65!onyx] at (59.0,18.7){+5,8};
 \fill[mustard] (36,14.25) rectangle ({36+0.55*4.24},15.75); \node[anchor=east,font=\fontsize{4.3}{4.6}\selectfont,text=spicy] at (59.0,15.0){+4,2};
 \fill[onyx!42] ({36-0.55*3.21},10.55) rectangle (36,12.05); \node[anchor=east,font=\fontsize{4.3}{4.6}\selectfont,text=muted] at (59.0,11.3){-3,2};
 \fill[spicy!65] ({36-0.55*12.73},6.85) rectangle (36,8.35); \node[anchor=east,font=\fontsize{4.3}{4.6}\selectfont,text=spicy] at (59.0,7.6){-12,7};
 \fill[spicy!42] ({36-0.55*10.91},3.15) rectangle (36,4.65); \node[anchor=east,font=\fontsize{4.3}{4.6}\selectfont,text=spicy!80!black] at (59.0,3.9){-10,9};
\end{tikzpicture}};
% Robustness callout
\node[anchor=north east,rounded corners=1pt,draw=rule,fill=white,inner sep=2pt,align=right,font=\sffamily\fontsize{5.35}{5.7}\selectfont] at (axis description cs:0.985,0.94){$\Delta$ decisión de misión: +20,9 p.p.\\excluyendo la ampliación dirigida: +22,0 p.p.};
\end{axis}
% right axis: MM3 mission-decision share
\begin{axis}[at={(patents.south west)},anchor=south west,width=0.95\linewidth,height=6.12cm,
 symbolic x coords={2012,2013,2014,2015,2016,2017,2018,2019,2020,2021,2022,2023,2024,2025},xmin=2012,xmax=2025,
 ymin=0,ymax=100,axis x line=none,axis y line*=right,ylabel={\sffamily\fontsize{6.1}{6.45}\selectfont MM3 decisión de misión (\%)},
 y label style={text=spicy,at={(axis description cs:1.075,0.5)},anchor=south},ytick={0,20,40,60,80,100},
 y tick label style={font=\sffamily\fontsize{5.4}{5.7}\selectfont,text=spicy,xshift=3pt},tick style={draw=spicy!70},axis line style={draw=spicy!70,line width=.35pt},clip=false]
\addplot+[spicy,line width=1.05pt,mark=o,mark size=1.35pt,mark options={fill=spicy,draw=spicy}] coordinates {(2018,53.3)(2019,50.0)(2020,41.7)(2021,35.3)(2022,37.9)(2023,40.0)(2024,47.2)(2025,58.8)};
\end{axis}
\end{tikzpicture}%
}}

\vspace{1.5mm}
\noindent\begin{minipage}[t]{0.43\linewidth}
{\sffamily\fontsize{8.05}{8.55}\selectfont\bfseries (b) Cobertura geográfica de solicitantes identificados}\par\vspace{0.4pt}
{\sffamily\fontsize{5.35}{5.65}\selectfont\color{muted}Titulares identificados: 2018--2022 ($n=74$) y 2023--2025 ($n=49$).}\par\vspace{0.7pt}
\resizebox{\linewidth}{!}{%
\begin{tikzpicture}[x=1mm,y=1mm,font=\sffamily\fontsize{4.75}{5.05}\selectfont]
\path[use as bounding box] (0,0) rectangle (73.5,49);
\def\xzero{22}\def\sc{0.67}
\foreach \v in {0,20,40,60}{\pgfmathsetmacro{\xx}{\xzero+\sc*\v}\draw[rule!55,line width=.2pt] (\xx,5)--(\xx,46);\node[anchor=north,text=muted,font=\fontsize{6.0}{6.4}\selectfont] at (\xx,4.3){\v};}
% legend
\fill[canary!80!onyx] (23,46.5) rectangle +(4,1.4); \node[anchor=west] at (28,47.2){2018--2022};
\fill[autumn] (44,46.5) rectangle +(4,1.4); \node[anchor=west] at (49,47.2){2023--2025};
% data macro: y/name/old/oldDirected/new/newDirected
\foreach \y/\name/\a/\ad/\b/\bd in {43/{China}/60.8/1.4/57.1/2.0,38.5/{EE.UU.}/2.7/2.7/14.3/14.3,34/{Japón}/5.4/1.4/10.2/6.1,29.5/{Suiza}/5.4/5.4/4.1/4.1,25/{Alemania}/4.1/0/6.1/0,20.5/{Francia}/4.1/4.1/4.1/4.1,16/{Canadá}/5.4/5.4/0/0,11.5/{Corea S.}/1.4/0/4.1/0,7/{Otros}/10.8/2.7/0/0}{
 \node[anchor=east,text=ink] at (16.9,\y){\name};
 \pgfmathsetmacro{\wa}{\sc*\a}\pgfmathsetmacro{\wad}{\sc*\ad}\pgfmathsetmacro{\wb}{\sc*\b}\pgfmathsetmacro{\wbd}{\sc*\bd}
 \fill[canary!80!onyx] (\xzero,\y+.25) rectangle ({\xzero+\wa},\y+1.55);
 \ifdim \ad pt>0pt \fill[pattern={Lines[angle=45,distance=1.35pt,line width=.22pt]},pattern color=onyx!70] ({\xzero+\wa-\wad},\y+.25) rectangle ({\xzero+\wa},\y+1.55);\fi
 \fill[autumn] (\xzero,\y-1.55) rectangle ({\xzero+\wb},\y-.25);
 \ifdim \bd pt>0pt \fill[pattern={Lines[angle=45,distance=1.35pt,line width=.22pt]},pattern color=onyx!70] ({\xzero+\wb-\wbd},\y-1.55) rectangle ({\xzero+\wb},\y-.25);\fi
 \ifdim \a pt>0pt \node[anchor=west,text=muted] at ({\xzero+\wa+.7},\y+.9){\pgfmathprintnumber[fixed,precision=1,use comma]{\a}};\fi
 \ifdim \b pt>0pt \node[anchor=west,text=spicy] at ({\xzero+\wb+.7},\y-.9){\pgfmathprintnumber[fixed,precision=1,use comma]{\b}};\fi
}
\draw[pattern={Lines[angle=45,distance=1.35pt,line width=.22pt]},pattern color=onyx!70,draw=onyx!60] (50.5,2.0) rectangle +(4,1.5);\node[anchor=west,text=muted] at (55.5,2.75){ampliación dirigida};
\node[anchor=north,text=muted,font=\fontsize{4.6}{4.9}\selectfont] at (41.5,0.8){\% dentro de cada periodo};
\end{tikzpicture}}
\end{minipage}\hfill
\begin{minipage}[t]{0.55\linewidth}
{\sffamily\fontsize{8.05}{8.55}\selectfont\bfseries (c) Orientación funcional por dominio de aplicación}\par\vspace{0.4pt}
{\sffamily\fontsize{5.35}{5.65}\selectfont\color{muted}Dominios explícitos ($n=75$): tres macrocapas funcionales.}\par\vspace{0.7pt}
\resizebox{\linewidth}{!}{%
\begin{tikzpicture}[x=1mm,y=1mm,font=\sffamily\fontsize{4.75}{5.05}\selectfont]
\path[use as bounding box] (0,0) rectangle (95,49);
\def\xstart{27}\def\barw{58}
\foreach \p in {0,25,50,75,100}{\pgfmathsetmacro{\xx}{\xstart+0.58*\p}\draw[rule!50,line width=.2pt] (\xx,5)--(\xx,45);\node[anchor=north,text=muted,font=\fontsize{6.0}{6.4}\selectfont] at (\xx,4.3){\p};}
\fill[autumn] (22,46.2) rectangle +(5,1.4);\node[anchor=west] at (28,46.9){misión};\fill[mustard] (41,46.2) rectangle +(5,1.4);\node[anchor=west] at (47,46.9){coordinación};\fill[onyx!42] (66,46.2) rectangle +(5,1.4);\node[anchor=west] at (72,46.9){física};
\foreach \y/\name/\m/\c/\p in {41.5/{Intralogística (39)}/69.2/17.9/12.8,35.5/{Manufactura (16)}/62.5/12.5/25.0,29.5/{Automoción (3)}/33.3/0/66.7,23.5/{Marítimo/obra (3)}/0/33.3/66.7,17.5/{Heavy-load (6)}/16.7/16.7/66.7,11.5/{Aeroespacial (8)}/12.5/0/87.5}{
 \node[anchor=east,text=ink] at (19.5,\y){\name};
 \pgfmathsetmacro{\wm}{0.58*\m}\pgfmathsetmacro{\wc}{0.58*\c}\pgfmathsetmacro{\wp}{0.58*\p}
 \fill[autumn] (\xstart,\y-1.55) rectangle ({\xstart+\wm},\y+1.55);
 \fill[mustard] ({\xstart+\wm},\y-1.55) rectangle ({\xstart+\wm+\wc},\y+1.55);
 \fill[onyx!42] ({\xstart+\wm+\wc},\y-1.55) rectangle ({\xstart+58},\y+1.55);
 \ifdim \m pt>18pt \node[font=\bfseries,text=white] at ({\xstart+0.5*\wm},\y){\pgfmathprintnumber[fixed,precision=0]{\m}};\fi
 \ifdim \c pt>18pt \node[font=\bfseries,text=onyx] at ({\xstart+\wm+0.5*\wc},\y){\pgfmathprintnumber[fixed,precision=0]{\c}};\fi
 \ifdim \p pt>18pt \node[font=\bfseries,text=white] at ({\xstart+\wm+\wc+0.5*\wp},\y){\pgfmathprintnumber[fixed,precision=0]{\p}};\fi
}
\node[anchor=north,text=muted,font=\fontsize{4.6}{4.9}\selectfont] at (50,1.0){\% del dominio};
\end{tikzpicture}}
\end{minipage}
  \caption[Evolución y estructura del corpus patentario curado.]{Evolución y
  estructura del corpus patentario curado. (a) Actividad anual, media móvil de
  tres años (MM3) de volumen y peso MM3 de decisión de misión; inset: cambio
  temático 2020--2022 $\rightarrow$ 2023--2025. (b) Cobertura geográfica de
  solicitantes identificados, con la ampliación dirigida por solicitante
  hachurada. (c) Orientación funcional por dominio de aplicación en tres
  macrocapas. Panel (b): top 8 más «Otros», con porcentajes normalizados dentro
  de cada periodo; el hachurado marca la ampliación dirigida. Panel (c): se
  excluyen 92 registros de uso general o multisector.}
  \label{fig:results-patents}
  \viusource{elaboración propia a partir del corpus patentario v3}
\end{figure}


El corpus patentario se desplaza hacia decisión de misión: media móvil de tres años del
37,9~\% al 58,8~\% entre el cierre de 2022 y el de 2025, es decir $+20{,}9$~puntos. El
signo y el orden de magnitud se mantienen en dos análisis de sensibilidad: excluyendo la
ampliación dirigida por solicitante, $+22{,}0$~puntos; colapsando los cuatro grupos de
títulos muy similares, $+20{,}0$. Mientras tanto, coordinación cinemática e interacción
física bajan (Figura~\ref{fig:results-patents}).

\begin{table}[!htbp]
  \centering
  \caption{Resultados de la revisión y objetivo específico al que condicionan.}
  \label{tab:results-review-synthesis}
  \fontsize{8.5}{10}\selectfont
  \renewcommand{\arraystretch}{1.06}
  \begin{tabularx}{\textwidth}{@{}>{\raggedright\arraybackslash}p{4.3cm}
      >{\raggedright\arraybackslash}p{1.5cm} Y@{}}
    \toprule
    Resultado de la revisión & Objetivo & Consecuencia para el diseño experimental \\
    \midrule
    Ningún trabajo del núcleo reúne certificación física, ejecución distribuida y reemplazo
      & OE2, OE3 & El certificado de \emph{wrench} y la reparación se evalúan como interfaces separadas y no se importa ninguna garantía de la literatura. \\
    La autoridad residual se desplaza a lo global en certificación y recuperación
      & OE1, OE5 & Toda dependencia global de la implementación evaluada se declara en la Tabla~\ref{tab:results-model-map} en vez de presentarse como distribuida. \\
    Doce trabajos con validación física frente a noventa y nueve en simulación
      & OE3, OE5 & Los resultados de simulación y los del demostrador se informan por separado, sin transferencia de conclusiones. \\
    Desplazamiento de +20,9 puntos hacia decisión de misión en el corpus patentario
      & OE1 & Muestra un desplazamiento reciente de la actividad de protección hacia la capa de decisión; no acredita adopción ni pertinencia industrial demostrada. \\
    Ningún precedente comercial integra selección local, certificación y recomposición
      & OE5 & El demostrador se plantea como prueba de compatibilidad funcional, no como comparación con un producto equivalente. \\
    \bottomrule
  \end{tabularx}
  \viuownsource
\end{table}

La Tabla~\ref{tab:results-review-synthesis} enlaza cada hallazgo con el objetivo que
condiciona. Ninguno demuestra que la arquitectura propuesta funcione: fijan qué debe
demostrarse y con qué comparadores, y excluyen importar garantías de trabajos con
supuestos distintos, presentar como distribuida una implementación con dependencias
globales, o equiparar simulación con validación física. Las subsecciones siguientes
evalúan las tres interfaces bajo esas
restricciones.
\zlabel{budget:review-end}
